\section{Integral formal series and bounded ordinary theta moments}
\label{ot-section}

This section proves the formal and functional-analytic parts of the
ordinary theta moment argument.  They are useful because bounded
ordinary derivatives alone do not define a bounded measure.  We give
the exact criterion, prove the descent from tensor powers, and prove
the integral functional equation used for the one-dimensional CM
theta function.  The geometric hypotheses needed to apply that
equation are stated explicitly at the end.

Fix a prime $p$.  Let $L$ be a complete nonarchimedean field containing
$\mathbb Q_p$, with its usual absolute value on $\mathbb Q_p$, and
write $\mathcal O=\{x\in L:|x|\le1\}$ and
$\mathfrak m=\{x\in L:|x|<1\}$.  A power series is called
\emph{integral} if all its coefficients belong to $\mathcal O$.
It is called \emph{bounded} if their absolute values have a finite
common bound.  No discreteness of the valuation is assumed.

\subsection{A complete measure criterion}

For $r\ge0$, put $X_r=\mathbb Z_p^r$.  For a multi-index
$n=(n_1,\ldots,n_r)\in\mathbb N^r$, set
\[
 B_n(x)=\prod_{j=1}^r\binom{x_j}{n_j},\qquad
 \binom Xk=\frac{X(X-1)\cdots(X-k+1)}{k!}.
 \tag{OT1}
\]
Empty products are one.  The polynomials $B_n$ take integral values
on $X_r$: this holds at nonnegative integers and then follows on
$\mathbb Z_p$ by continuity and density.
Let $S_jf(x)=f(x+e_j)$ and $\Delta_j=S_j-1$.

\begin{theorem}[Multivariable Mahler expansion]\label{ot-mahler}
Every continuous $f:X_r\to L$ has a unique uniformly convergent
expansion
\[
 f(x)=\sum_{n\in\mathbb N^r}a_n(f)B_n(x),\qquad
 a_n(f)=(\Delta_1^{n_1}\cdots\Delta_r^{n_r}f)(0).
 \tag{OT2}
\]
The coefficients tend to zero outside finite subsets of
$\mathbb N^r$, and
\[
                  \|f\|_\infty=\sup_n|a_n(f)|.
 \tag{OT3}
\]
When $r=0$, these assertions mean the single constant coefficient.
\end{theorem}
\begin{proof}
The ultrametric inequality gives $\|\Delta_j\|\le1$, hence
$|a_n(f)|\le\|f\|_\infty$.  We first prove their convergence
to zero.  A continuous function on the compact space $X_r$ can be
approximated uniformly, to any accuracy, by a function constant on
residue classes modulo $p^h$.  Indeed, uniform continuity gives
such an $h$, and one chooses one value on each of the finitely many
classes.

Let $g$ be such a function and put $N=p^h$.  On these functions
$S_j^N=1$.  In the ring $\mathbb Z[S]/(S^N-1)$ all coefficients
of $(S-1)^N$ are divisible by $p$: reduction modulo $p$ gives
$(S-1)^{p^h}=S^{p^h}-1=0$.  It follows that
\[
 \|\Delta_j^{kN}g\|_\infty\le |p|^k\|g\|_\infty.
 \tag{OT4}
\]
The same estimate holds with any additional differences in any
coordinates.  Thus $a_n(g)$ tends to zero whenever at least one
coordinate of $n$ tends to infinity, uniformly in the others.
Since $|a_n(f)-a_n(g)|\le\|f-g\|_\infty$, the asserted
cofinite convergence follows for $f$.

The series in (OT2) therefore converges uniformly: its tails have
norm at most the largest omitted coefficient, because
$|B_n(x)|\le1$.  Its sum is continuous.  For
$x\in\mathbb N^r$, the series is finite, and the ordinary Newton
identity gives its value as $f(x)$.  In one variable this identity
is the binomial inversion
\[
 f(m)=\sum_{n=0}^m\binom mn
          \sum_{j=0}^n(-1)^{n-j}\binom nj f(j).
\]
The coefficient of $f(j)$ is
$\binom mj\sum_{k=0}^{m-j}(-1)^k\binom{m-j}{k}$, which is
one for $j=m$ and zero otherwise.  Applying this in each coordinate
proves the multivariable identity.  Density of $\mathbb N^r$ in
$X_r$ proves (OT2) everywhere.  It also gives
$\|f\|_\infty\le\sup_n|a_n(f)|$, establishing (OT3).
Finally $\Delta_j\binom{x_j}{k}=\binom{x_j}{k-1}$ implies
$a_n(B_m)=1$ if $n=m$ and zero otherwise.  Applying the bounded
functional $a_n$ to any uniformly convergent expansion proves
uniqueness.
\end{proof}

By an $L$-valued bounded measure on $X_r$ we mean a continuous
$L$-linear functional on $C(X_r,L)$, with its supremum norm.  This
definition avoids any assumption of countable additivity in a
real-valued measure space.

\begin{corollary}[Exact moment criterion]\label{ot-measure}
A family $(c_n)_{n\in\mathbb N^r}$ is of the form
$c_n=\mu(B_n)$ for a bounded measure if and only if it is bounded.
The measure is unique, is given by
\[
             \mu(f)=\sum_n a_n(f)c_n,
             \qquad \|\mu\|=\sup_n|c_n|,
 \tag{OT5}
\]
and has norm at most one exactly when every $c_n$ is integral.
\end{corollary}
\begin{proof}
Boundedness of the $c_n$ and the cofinite convergence of $a_n(f)$
give convergence of (OT5) in the complete field $L$.  The
ultrametric inequality and (OT3) show that it is a linear
functional of norm at most $\sup|c_n|$.  Since
$\|B_n\|_\infty=1$ and $a_m(B_n)=\delta_{mn}$, the reverse
inequality holds.  Conversely a bounded measure has
$|\mu(B_n)|\le\|\mu\|$.  Applying it to (OT2) proves
uniqueness.
\end{proof}

Let $g\in L[[z_1,\ldots,z_r]]$ and write it in exponential
generating-function form
\[
 g(z)=\sum_{\alpha\in\mathbb N^r}
           M_\alpha\frac{z^\alpha}{\alpha!}.
 \tag{OT6}
\]
Here $\alpha!=\prod_j\alpha_j!$; it is nonzero in $L$.
Define the linear functional $\Lambda_g$ on $L[X_1,\ldots,X_r]$
by $\Lambda_g(X^\alpha)=M_\alpha$.

\begin{proposition}\label{ot-logarithm}
The coefficients of the formal series
\[
       g(\log(1+T_1),\ldots,\log(1+T_r))
       =\sum_n c_nT^n
 \tag{OT7}
\]
are exactly $c_n=\Lambda_g(B_n)$.  Consequently $\Lambda_g$
extends to a measure on $X_r$ of norm at most one if and only if
the series (OT7) is integral.
\end{proposition}
\begin{proof}
The formal identity
\[
 \exp(X\log(1+T))=\sum_{n\ge0}\binom Xn T^n
 \tag{OT8}
\]
holds in $L[X][[T]]$.  One proof differentiates in $T$: both sides
have constant term one and satisfy
$(1+T)\partial_T F=XF$, whose coefficients determine the series
uniquely.  Multiply this identity in the $r$ coordinates and apply
$\Lambda_g$ coefficient by coefficient.  Each coefficient uses
only finitely many monomials, so the operation is purely formal
and gives (OT7).  Corollary~\ref{ot-measure} proves the criterion.
Polynomial uniqueness follows also from the fact that the $B_n$
form a basis, since their leading monomials are $X^n/n!$.
\end{proof}

For a complex analytic function $G$ near zero in a finite-dimensional
complex vector space, choose vectors $v_1,\ldots,v_r$, an abstract
field embedding of its Taylor coefficients into $L$, and a period
$\omega\in L$.  Apply the proposition to the transported Taylor
series of $G(\omega\sum_j z_jv_j)$; equivalently set
$M_\alpha=\omega^{|\alpha|}\iota(D_{v_1}^{\alpha_1}
\cdots D_{v_r}^{\alpha_r}G(0))$.
Only algebraic operations on Taylor coefficients are involved; the
field embedding is not asserted to be continuous for the complex
topology.

\subsection{Descent from tensor powers}

\begin{theorem}\label{ot-tensor-root}
Let $G$ be as above, let $\Lambda$ be an additive subgroup of its
vector space, and let $h\ge1$.  Define the function in independent
blocks
\[
             G^{\otimes h}(z_1,\ldots,z_h)
                      =\prod_{j=1}^hG(z_j).
 \tag{OT9}
\]
If, for every finite family of directions in $\Lambda^h$, its
period-scaled Taylor functional is a measure of norm at most $C$,
then the corresponding functional of $G$ in every finite family
of directions in $\Lambda$ is a measure of norm at most $C^{1/h}$.
Here $C\ge0$.  In particular norm at most one descends from the
fourth tensor power, without any assumption on $G(0)$.
\end{theorem}
\begin{proof}
Fix $v_1,\ldots,v_r\in\Lambda$ and let $c_n$ be the coefficients
in (OT7) for these directions.  Use in the tensor product the
$hr$ directions obtained by placing $v_i$ in block $j$ and zero
in the other blocks.  Its logarithm-substituted Taylor series is
the product of $h$ copies of (OT7) in disjoint variable sets.
The coefficient indexed by $(n,\ldots,n)$ is therefore $c_n^h$.
The assumed norm bound and Corollary~\ref{ot-measure} give
$|c_n|^h\le C$.  Thus $|c_n|\le C^{1/h}$, and the same
corollary gives the claimed measure.  For $C=0$ every $c_n$ is
zero; this case is included.  For an empty direction family the
argument is simply the same inequality for the constant term.
\end{proof}

The independent blocks in this theorem matter.  A pointwise power
$G(z)^h$ combines coefficients by convolution and does not provide
the diagonal coefficient $c_n^h$.  The submitted fourth-power
argument uses the tensor product (OT9), so the theorem applies
with precisely its intended meaning.

\subsection{An integral functional equation}

\begin{theorem}\label{ot-functional-equation}
Let $\phi=(\phi_1,\ldots,\phi_r)\in
\mathcal O[[T_1,\ldots,T_r]]^r$ have zero constant term and
linear coefficient matrix in $\operatorname{GL}_r(\mathcal O)$.
Let $m\ge0$ be an integer with $|m|<1$, and let
$A,B\in\mathcal O[[T_1,\ldots,T_r]]$ have constant term one.
There exists a unique $f\in L[[T_1,\ldots,T_r]]$, with constant
term one, satisfying
\[
                     f(\phi(T))A(T)=f(T)^mB(T).
 \tag{OT10}
\]
This solution is integral.  If $m\ge2$ and
$0\ne\rho\in L[[T_1,\ldots,T_r]]$ satisfies
$\rho(\phi)=\rho^m$, then $\rho$ is a constant $c$ with
$c^{m-1}=1$.
\end{theorem}
\begin{proof}
Write the linear part of $\phi$ as $UT$.  Suppose the homogeneous
parts of $f$ below degree $n$ have been constructed.  The degree
$n$ part $f_n$ occurs on the left of (OT10) as $f_n(UT)$ and on
the right as $mf_n(T)$.  All other degree $n$ terms involve only
the already constructed parts.  Thus
\[
                       (S_n-m)f_n=R_n,
 \tag{OT11}
\]
where $S_n:P(T)\mapsto P(UT)$ is an automorphism of the finite
free $\mathcal O$-module of homogeneous polynomials of degree
$n$.  Its inverse substitutes $U^{-1}T$, so both $S_n$ and
$S_n^{-1}$ have integral entries.  The matrix
$1-mS_n^{-1}$ reduces to the identity modulo $\mathfrak m$;
its determinant is a unit of $\mathcal O$, and its inverse is
integral by the adjugate formula.  Hence $S_n-m$ is invertible
over $\mathcal O$.  If the lower parts of $f$ are integral, then
$R_n$ is integral, so (OT11) determines a unique integral $f_n$.
Starting with $f_0=1$ proves existence, integrality and uniqueness
degree by degree.  The same invertibility proves uniqueness among
all $L$-valued formal solutions, even if their coefficients were
not initially assumed bounded.  The case $r=0$ is the constant
identity $1=1$.

For the last assertion, the least degree of a nonzero homogeneous
part of $\rho$ is unchanged by substitution of $\phi$, since
$U$ is invertible.  Equation $\rho(\phi)=\rho^m$ forces this
degree to be zero.  Its constant $c$ is nonzero and satisfies
$c=c^m$.  Then $\rho/c$ satisfies (OT10) with $A=B=1$ and
constant term one.  Uniqueness makes it the constant one.
\end{proof}

In one variable the recurrence is especially transparent: if
$\phi(T)=aT+\cdots$, the coefficient of $f_n$ in (OT11) is
$a^n-m$, whose absolute value is one.  This proves exactly the
integrality step used in the ordinary CM theta construction and
also shows why the unit linear coefficient is essential.

\subsection{Bounded elliptic quotients}

The next elementary observation supplies the bounded-series step
in that construction.  We isolate the geometric inputs so that
integrality of a formal chart is not tacitly assumed.

\begin{lemma}\label{ot-bounded-unit}
Suppose $F,G\in L[[T]]$ are bounded, $FG=1$ and $F(0)=1$.
Then both $F$ and $G$ are integral.
\end{lemma}
\begin{proof}
For $0<r<1$ define $\|F\|_r=\max_n|F_n|r^n$; the maximum
exists because $F$ is bounded.  The same applies to $G$.
This norm is multiplicative.  Indeed, choose the smallest indices
$i,j$ where the two maxima occur.  In the coefficient of
$T^{i+j}$ in $FG$, every other summand has smaller weighted
absolute value: if its first index is less than $i$ it is already
smaller, and otherwise its second index is less than $j$.
The ultrametric inequality then gives equality for this
coefficient, while the reverse inequality for all coefficients
is immediate.  Thus $\|F\|_r\|G\|_r=1$.
Both norms are at least one, from their constant terms, hence
both are one.  Letting $r$ tend to one gives the assertion for
each coefficient.
\end{proof}

\begin{proposition}\label{ot-elliptic-bounded}
Let $E/L$ be a smooth Weierstrass elliptic curve, with origin $O$,
and suppose a formal parameter $T$ at $O$ gives expansions
\[
 x=T^{-2}X(T),\qquad y=T^{-3}Y(T),\qquad
 X,Y\in\mathcal O[[T]],\quad X(0)\in\mathcal O^\times.
 \tag{OT12}
\]
Let $R$ be a rational function on $E$, regular at $O$.
Assume every pole $P$ of $R$ is affine and satisfies $|x(P)|\le1$,
after passage to a complete algebraically closed extension if needed.
Then the expansion of $R$ in $T$ is bounded.
If $R(O)=1$ and the same pole condition holds for $R^{-1}$, then
the expansion of $R$ and its inverse are integral.
\end{proposition}
\begin{proof}
For each pole choose a sufficiently large multiplicity and form a
monic polynomial $D(X)$ whose roots are the finitely many $x(P)$.
Then $D(x)R$ has no affine pole, and $D$ has integral coefficients.
The affine coordinate ring of the smooth Weierstrass curve is
$L[x,y]$ modulo its Weierstrass equation.  A rational function
regular at every affine point belongs to this ring.  To recall
the local justification, smoothness makes each local ring a
discrete valuation ring, so absence of a pole means membership in
that local ring.  If an element of the fraction field belongs to
all those local rings but not to the affine ring $A$, its
denominator ideal $\{a\in A:af\in A\}$ is proper.  A maximal
ideal containing it contradicts membership in the corresponding
localization.  Thus
\[
                        D(x)R=A(x)+yB(x)
 \tag{OT13}
\]
for polynomials $A,B$, since the Weierstrass equation is monic of
degree two in $y$.

Put $d=\deg D$.  From (OT12),
$D(x)=T^{-2d}S(T)$, where $S$ is integral and
$S(0)=X(0)^d$ is a unit.  Its inverse is integral by the usual
coefficient recursion.  Multiplying the right side of (OT13)
by a sufficiently large power $T^k$ produces a bounded series:
it is a finite sum of scalar multiples of products of $X$ and
$Y$.  Therefore the Laurent expansion
$R=T^{2d-k}(T^k(A(x)+yB(x)))S^{-1}$ has bounded coefficients.
Regularity at $O$ removes its negative coefficients, yielding the
claimed power series.  Applying the same argument to $R^{-1}$
and Lemma~\ref{ot-bounded-unit} proves the last assertion.
All arguments may be made after the indicated extension; bounds
and integrality of the original coefficients then descend.
\end{proof}

\subsection{The precise one-dimensional CM consequence}

Here is a sufficient set of ordinary CM data, stated as conventional
mathematics.  Let $k$ be an imaginary quadratic field, let $p$ be
odd and split in $k$, and fix the embedding
$\mathcal O_k\hookrightarrow\mathbb Z_p$ associated with one
of its primes.  Suppose a complex elliptic curve with CM by
$\mathcal O_k$ has been transported to a good integral Weierstrass
model over $L$.  Assume its formal group has an integral
multiplicative coordinate $T$: its parameter at the origin is an
integral series in $T$ with unit linear coefficient, and the CM
endomorphism $\alpha$ acts as
\[
                         [\alpha](T)=(1+T)^{a}-1,
             \qquad a=\iota(\alpha)\in\mathbb Z_p.
 \tag{OT14}
\]
Assume the CM endomorphisms preserve the formal subgroup, and that
this subgroup consists precisely of $O$ and the affine points with
$|x|>1$.  Suppose a normalized entire theta function $\theta$
has $\theta(0)=1$, zeros at the translates of a nonzero point of
order two, and that
\[
              R_\alpha(z)=\frac{\theta(\alpha z)}
                                     {\theta(z)^{N(\alpha)}}
 \tag{OT15}
\]
is an elliptic rational function for each nonzero
$\alpha\in\mathcal O_k$ used below.  Finally suppose the
transported Taylor series in the formal coordinate is
\[
               f(T)=\widehat\theta(\omega\log(1+T))
 \tag{OT16}
\]
with the same period $\omega$ as in the differentiation
directions; its constant term is one.  These are hypotheses on
the actual model and normalization, not consequences of bounded
formal algebra.

\begin{corollary}\label{ot-CM-base}
Under these hypotheses, $f$ is integral.  Consequently the Taylor
functional of $\theta$ in the period-scaled one-dimensional
direction is a measure of norm at most one on $\mathbb Z_p$.
\end{corollary}
\begin{proof}
The Chinese remainder theorem gives
$\alpha\in\mathcal O_k$ which is one modulo the chosen prime
over $p$ and zero modulo its conjugate.  Thus $|a|=1$ and
$m=N_{k/\mathbb Q}(\alpha)$ is a positive integer divisible by
$p$.  In particular $m\ge2$ and $|m|<1$.

A good integral formal group has no nonzero prime-to-$p$ torsion.
In fact its multiplication-by-$n$ series is
$nU+U^2H(U)$ with $H$ integral; when $p\nmid n$ and
$0<|U|<1$, its absolute value is $|U|$, so it cannot vanish.
Since $p$ is odd, the theta zero of order two does not lie in the
formal subgroup.  No preimage of this zero under a CM
endomorphism lies there either, because that subgroup is
preserved.  The zeros and poles of $R_\alpha$ therefore avoid
the formal subgroup.  Also $R_\alpha(O)=1$.
Proposition~\ref{ot-elliptic-bounded} makes the expansion
$\epsilon(T)$ of $R_\alpha$ integral, with integral inverse.

The coefficients of $(1+T)^a-1$ are integral, by the
integer-valued binomial polynomials and density of $\mathbb N$
in $\mathbb Z_p$; its linear coefficient is the unit $a$.
Equations (OT14)--(OT16), as formal identities, give
\[
              f((1+T)^a-1)=f(T)^m\epsilon(T).
 \tag{OT17}
\]
Theorem~\ref{ot-functional-equation}, with $A=1$ and
$B=\epsilon$, proves that $f$ is integral.
Proposition~\ref{ot-logarithm} now gives the asserted measure.
\end{proof}

The extra hypotheses are necessary to identify this corollary with
the theta function used in a global construction.  In particular
one must prove existence and normalization of the good CM model
and multiplicative chart, that its period is a unit in the chosen
normalization, integrality for the required prime-to-$p$ torsion
translates and their finite linear combinations, and the passage
from the one-dimensional theta function to the higher-dimensional
lattice theta function.  The tensor descent in
Theorem~\ref{ot-tensor-root} closes the final extraction from the
fourth tensor power; it does not supply that tensor power's
integrality.  The original proof's treatment of $p=2$ requires
additional information on its specific point of order two, and
is not inferred from Corollary~\ref{ot-CM-base}.

\subsection{Checks against stronger false assertions}

\begin{enumerate}
\item Bounded power moments are insufficient.  The functional
$P\mapsto P'(0)$ has power moments zero except for the moment of
$X$, which is one, but takes $\binom X{p^a}$ to
$(-1)^{p^a-1}/p^a$ for every $a\ge1$.  These values have
unbounded absolute value on polynomials of supremum norm one,
so there is no bounded extension.  The binomial, rather than power, moments in
Corollary~\ref{ot-measure} are essential.
\item The invertible linear part in
Theorem~\ref{ot-functional-equation} cannot be discarded.  In
$\mathbb Q_p[[T]]$, put
$f(T)=\exp(T/p)$, $\phi(T)=pT$ and $m=p$.  Then
$f(\phi)=f^p$, with $A=B=1$, although $f$ is not integral.
\item The contraction condition $|m|<1$ is also essential.
For $\phi(T)=T$ and $m=1$, every series with constant term one
satisfies the equation with $A=B=1$, including $1+p^{-1}T$.
\item Boundedness of a series does not imply boundedness of its
reciprocal: $1+p^{-1}T$ is bounded, but its reciprocal has
coefficients $(-p^{-1})^n$.  Both bounds are used in
Lemma~\ref{ot-bounded-unit}.
\item At $p=2$ a formal group can have nonzero torsion of order
two: in the multiplicative formal group the parameter $T=-2$
has $|T|_2<1$ and $(1+T)^2-1=0$.  Thus the last proof's
prime-to-$p$ torsion argument genuinely requires odd $p$.
\end{enumerate}


